\newpage
% chapter 0
\chapter*{Prefazione}
\addcontentsline{toc}{chapter}{Prefazione}

Questo documento è stato realizzato per fornire una dispensa a coloro che si stanno approcciando allo studio delle reti di calcolatori.
Si tratta di una trascrizione dei concetti fondamentali trattati nel corso di \textit{Reti di Calcolatori} tenuto dal Prof. \textit{Leonardo Maccari} presso l'Università Ca' Foscari di Venezia nell'anno accademico 2025/2026.

\vspace{0.5 em}

Gli appunti sono tratti da lezioni frontali, slide, libri di testo del corso e analisi personali, che spero siano corrette e che possano essere utili per approfondire i temi trattati e chiarire eventuali dubbi.
Le immagini e gli schemi presenti nel documento sono stati creati personalmente o sono tratti dalle slide del corso e dai libri di testo consigliati. Tali contenuti sono utilizzati esclusivamente a scopo didattico e illustrativo, restando di proprietà dei rispettivi autori.

\vspace{0.5 em}

Inoltre, ci tengo a precisare che il documento presentato non sostituisce lo studio dei testi consigliati, né tantomeno la frequenza delle lezioni, che permettono di comprendere appieno gli argomenti affrontati durante il corso.
Spero che le spiegazioni siano chiare e che non siano presenti errori dovuti a interpretazioni personali, anche se non garantisco nulla.

\vspace{0.5 em}

Questo ebook è disponibile anche al link \extlink{https://github.com/giovanni-quacchia/Computer-Networks-Notes}{Github} e può essere condiviso liberamente per usi non commerciali, purché non vengano apportate modifiche al contenuto originale e venga sempre citata la fonte.